%===================================================================
% Chapter 6: Future Work
%===================================================================
\chapter{Future Work}
\label{chap:future}

The systems in this dissertation answer the question they were built
for: SQL can be rewritten, correctly and profitably, beyond any fixed
set of rules. The answer opens questions of its own. This chapter
outlines five directions, in each case starting from something this
dissertation already demonstrates and asking what lies one step
further.

\section{From Rewrites to Rules and Shared Optimization Knowledge}
\label{sec:fw:rules}

Rule-free rewriters and rule-based rewriters need not remain
competitors: the former can supply the latter. Every rewrite this
dissertation's systems produce is a verified, measured example of a
transformation that existing rule sets lack, and such examples are
precisely the raw material from which new rules can be distilled. Two
concrete paths follow from Chapters~3 and~4. SlabCity's
dataflow-guided synthesis can be pointed at rule discovery directly:
a specific transformation found for one query can be generalized into
a template pair with an associated condition, and verified once as a
rule in the style of WeTune~\cite{wetune}. GenRewrite's repository
suggests a complementary path: each NLR2 accumulates multiple
verified query pairs over time, and aggregating those pairs makes it
possible to extract the common structural transformation together
with explicit preconditions, so that the resulting rule generalizes
beyond the incidental syntax of any single example. A rule mined this
way arrives with evidence: the queries it improved and the speedups
it delivered. NLR2s add one more possibility that classical rules
never had. Because they are stated in natural language and contain no
schema or data from the workload that produced them, they can be
shared across teams and organizations, raising the prospect of
community repositories of optimization knowledge that grow the way
open-source code does.

\section{Optimizing AI-Generated SQL and Learning to Rewrite}
\label{sec:fw:ai-sql}

The workload is changing under this field's feet. An ever larger
share of SQL is produced by AI, through assistants and natural
language interfaces, and AI-generated SQL fails differently from
human SQL: it is written to capture intent, by a generator that
cannot see the data's statistics or the engine's cost behavior, so
its inefficiencies are systematic and repeated at scale. Three
research efforts follow, in increasing ambition.

The first is to understand the new workload. Classical benchmarks
encode human idioms, so the field lacks a benchmark of AI-generated
SQL; the generated query variants in Chapter~4's evaluation are a
seed, and expanding them into a curated benchmark, paired with a
taxonomy of the inefficiency patterns generators actually produce,
would give rewriting research its proper target and quantify how far
techniques tuned on human queries carry over.

The second is to make rewriting knowledge flow backward into
generation. The NLR2 repository is exactly the knowledge a generator
lacks: supplying relevant NLR2s at generation time, or screening
candidate generations against known anti-patterns before they enter
a codebase, would prevent inefficient SQL rather than cure it.
Prevention will not eliminate the cure, since cost depends on data
and engine characteristics that generation time cannot see, but
measuring how much of the rewriting gap prevention closes is itself
an open question.

The third is to learn the rewriter. The most capable AI models
today are trained by reinforcement learning on tasks with
automatically checkable success, and query rewriting is such a task:
equivalence can be checked and speedup can be measured, so this
dissertation's checking machinery is, from a training perspective, a
reward function. The verified rewrite pairs GenRewrite accumulates
are training data that no one else has. A concrete program starts
with supervised fine-tuning on those pairs, continues with
reinforcement learning in which a candidate rewrite is rewarded only
when the tester and verifier accept it and the measured or estimated
cost improves, and aims at a small, specialized rewriting model that
has internalized what GenRewrite currently retrieves through
prompting. Such a model would be cheaper per query than a
general-purpose model, deployable privately beside the warehouse,
and, unlike a rule set, would keep improving as its own accepted
rewrites become further training signal.

\section{The Continuously Optimized Pipeline}
\label{sec:fw:pipelines}

Chapter~5 optimizes a pipeline as it exists at one moment, but
pipelines evolve for years, and evolution interacts with
optimization in ways a one-time pass cannot address. The core
technical problem is dependence tracking. A refactoring that merges
models or moves logic across nodes is correct relative to the
producer definitions it read at optimization time; a later edit to
one of those definitions can silently invalidate it. A continuous
optimizer must therefore record, for every accepted rewrite, the
definitions it depends on, detect when an edit touches them, and
re-verify only the affected subgraph, using the same data-backed
checks as the original decision. The workflow question is equally
concrete: the natural insertion point is code review, where the
optimizer analyzes a proposed change to the pipeline, re-validates
the optimizations it affects, and proposes new refactorings as
ordinary review comments with diffs attached, so optimization
becomes part of how the pipeline is developed rather than a
maintenance event. Beyond correctness, a continuous optimizer must
manage change itself: refresh frequencies and data volumes drift, so
materialization and scheduling decisions that were right last
quarter decay, and the system needs stability rules that re-tune
when the projected savings justify the churn rather than
oscillating with every fluctuation. Solving these problems would
turn DAGSmith from a tool one runs into a steward that accompanies
a pipeline through its life, which is the form in which pipeline
optimization would actually be adopted.

\section{Query Optimization as a Cloud Service}
\label{sec:fw:service}

The economics that motivate this dissertation point at a delivery
model for it. In cloud warehouses, optimization converts directly
into money saved, which makes rewriting natural to offer as a
service that watches an organization's recurring workload, proposes
verified rewrites and pipeline refactorings, and is judged, perhaps
even priced, by the compute it saves. Turning that sketch into
reality raises research questions rather than merely engineering
ones. The first is attribution: charging for savings requires
measuring them credibly, and since costs shift with data growth and
workload drift, the service needs counterfactual estimates of what
the unoptimized workload would have cost, validated by controlled
comparisons on recurring runs. The second is trust and autonomy:
organizations will not hand write access to their pipelines to an
external system, so the service must span a spectrum from advisory
mode to auto-apply with verified rollback, with the equivalence
machinery of this dissertation serving as the guardrail that makes
higher autonomy defensible. The third is what multi-tenancy does to
knowledge: a service that optimizes many customers accumulates
NLR2-style knowledge across all of them, and because NLR2s carry no
customer schema or data, each tenant's successful rewrites can
safely make every other tenant's optimization better, giving the
service a compounding advantage no single-tenant tool can match, and
giving a new customer useful optimization from day one. There is
also an ecosystem question worth naming: the warehouse vendor's
revenue grows with compute consumed while the service's value grows
with compute saved, and where optimization ultimately lives, inside
the warehouse, beside it, or independent of it, will be decided by
that tension.


\section{Verification in the AI Era}
\label{sec:fw:verification}

The guarantees in this dissertation are bounded by their verifiers:
full verification reaches 17\% of Chapter~3's rewrites, the most
stubborn cost in Chapter~4 is the human expert who must sign off where
the verifier cannot, and Chapter~5's pipeline equivalence rests on
data-backed evidence rather than proof. Across all three systems the
pattern is the same: proposing rewrites has become cheap and
plentiful, while proving them correct remains scarce and expensive.
Verification is the bottleneck of automated query rewriting.

There is concrete reason to believe this bottleneck will move. AI
systems reached gold-medal standard at the International Mathematical
Olympiad in 2025~\cite{luong2025imogold}, and within the past year
LLM-based theorem provers have produced formal, machine-checked proofs
at scale and contributed to the resolution of open research problems,
including longstanding open problems posed by
Erd\H{o}s~\cite{tsoukalas2026proofsearch}. Proving query equivalence
is mathematics of a far shallower kind: when the verifiers in this
dissertation fail, it is rarely because the underlying theorem is
deep, but because their mechanized semantics do not cover the
operators involved or their decision procedures time out on formulas a
human could dispatch with a short argument. The barrier is breadth and
engineering, and those are what learned systems scale.

The architecture of AI-era verification is already visible in these
mathematical systems, and it follows the same principle this
dissertation applied to rewriting: the model proposes, and something
trustworthy disposes. Equivalence under a mechanized SQL semantics, in
the style pioneered by \cosette~\cite{chu2017cosette}, is stated as a
theorem; a language model searches for a proof; and a small trusted
kernel, a Lean-style proof checker or a solver certificate checker,
validates the result. A hallucinated proof is then not a wrong verdict
but a rejected artifact that fails to check, so soundness resides
entirely in the kernel and the model contributes only search, which is
precisely what models have become good at. The same models attack the
coverage gap from the other side: autoformalization, which already
translates mathematical statements into formal ones with high
accuracy~\cite{gao2025herald}, can draft the mechanized semantics of
SQL operators and dialect variations from their documentation and
reference behavior, each audited once by a human and reused forever,
so that verifier coverage grows at the pace of a review queue rather
than a research career.

Equivalence for the full language remains undecidable, and NULL
semantics, aggregation, and window functions will stay the hard cases;
no breakthrough repeals that. But if machine-found, machine-checked
equivalence proofs become the common case, the consequences reach
every chapter of this dissertation: Chapter~3's label taxonomy
collapses upward as fully verified becomes the default, the human
expert of Chapter~4 is consulted about intent rather than equivalence,
and the reinforcement-learning program of \cref{sec:fw:ai-sql}
acquires an exact reward signal in place of an approximate one. Given
the current rate of progress in AI mathematics, a breakthrough within
the next few years would not be surprising, and a field whose central
bottleneck sits directly in its path should prepare for it.
